\section{Conclusion}

We built the multimodal core of BaiZe on the frozen BaiZe language backbone~\cite{baiZe-llm}: a vision encoder pre-trained from scratch, aligned to the backbone through a lightweight projector, and adapted to the semiconductor domain via preference-based reinforcement learning.
% 中文：我们在冻结的 BaiZe 语言骨干 \cite{baiZe-llm} 之上构建了 BaiZe 的多模态内核：从零预训练的视觉编码器，经轻量投影层对齐到骨干，再通过基于偏好的强化学习适配到半导体领域。

Through a controlled encoder comparison, we found that a larger ViT-L/14 tower ($68.2\%$ linear-probe, $+3.9$\,pp over ViT-B/16) does \emph{not} translate its advantage downstream: the aligned MLLM averages $26.1\%$ over 11 benchmarks, indistinguishable from the smaller tower. We attribute the gap to the $224$-px single-tile setting and the alignment architecture rather than to tower capacity.
% 中文：通过受控的编码器对比，我们发现更大的 ViT-L/14 塔（线性探测 $68.2\%$，比 ViT-B/16 高 $+3.9$\,pp）并未把优势迁移到下游：对齐后的 MLLM 在 11 个基准上平均 $26.1\%$，与小塔不可分辨。我们把这一落差归因于 224 像素单块设定与对齐架构，而非塔容量。

Stages~(ii) and~(iii) are evaluated across 12 general multimodal benchmarks, 4 domain benchmarks, and the multimodal subset of an EDA flow benchmark. Stage~(iii) is reported here as a CPU-only pilot (Section~\ref{sec:domain}): it yields one controlled, statistically significant gain on a waveform type-classification probe ($+16$\,pp over SFT, $p{=}0.0009$, cheat rate $0$) alongside a sequence of honestly reported non-significant or negative iterations, and an audit that bounds the ceiling of an image channel on the current benchmark. Scaling the significant pairing to GPU-scale RL, and the cross-line VQA-tool re-evaluation of the waveview cases, remain open and are stated as such rather than estimated.
% 中文：第 (ii) 与 (iii) 阶段在 12 个通用多模态基准、4 个领域基准与 EDA 流程基准的多模态子集上评测。第 (iii) 阶段在此作为仅 CPU 试点报告（第 \ref{sec:domain} 节）：它在波形类型分类探针上给出一个受控统计显著增益（较 SFT +16pp，$p=0.0009$，cheat rate 0），同时如实报告一连串非显著或负的迭代，以及一份界定「图像通道在当前基准上天花板」的审计。把显著组合扩展到 GPU 规模 RL、以及对 waveview 用例的跨线 VQA 工具再评测，仍属开放，据实陈述而非估算。

The immediate next step is closing the encoder-to-downstream transfer gap---through higher-resolution multi-tile input and domain-specific alignment---toward a deployable multimodal core for visual EDA tasks.
% 中文：下一步工作是弥合「编码器到下游」的迁移落差——通过更高分辨率的多块输入与领域专项对齐——目标是得到一个可用于视觉 EDA 任务的可部署多模态内核。
